\section{English Wikipedia as a multi-document summarization dataset}
Wikipedia, being an encyclopedia, can be viewed as a collection of
summaries on various topics given by their title,
e.g. "Canada" or "Machine Learning". The source material to be 
summarized can be viewed as all reputable documents on the Web or
books; however, to make the problem more tractable we consider the following
subsets of all documents, $D$: 

\begin{enumerate}
    \item Cited sources:
A Wikipedia article that conforms to the style
guidelines should be well-supported by citations found in the
\textit{References} section of Wikipedia articles.
For each article, $a_i$, we extract all text without
markup from crawlable citation documents, $C_i \subset D$, to use as input
to our method.
\item Web Search results:
To expand the collection of reference documents, we crawl the search results
from the Google search engine, using the article section titles as queries.
For each query, we collect 10 result pages. From this collection we remove
the Wikipedia article itself, which is often among the top results.
We also remove "clones", which are detected
when there is a high-level of unigram overlap with the article (details provided in \ref{clones}).
We denote these refined search results for an article, $a_i$, as $S_i \subset D$.
Similar to $C_i$, we extract only the text to use as input.
\end{enumerate}

\begin{table}[t]
\caption{Percentiles for different aspects of WikiSum dataset. Size is in number of words.}
\begin{center}
\begin{tabular}{rrrrrrr}
\multicolumn{1}{r}{\bf Percentile}&\multicolumn{1}{r}{\bf 20}&\multicolumn{1}{r}{\bf 40}&\multicolumn{1}{r}{\bf 50}&\multicolumn{1}{r}{\bf 60}&\multicolumn{1}{r}{\bf 80}&\multicolumn{1}{r}{\bf 100}\\ \hline \\
Lead Size	&37	&62	&78	&98	&166	&10,034 \\
Num Citations	&1	&2	&2	&3	&5	&1,029 \\
Citations Size	&562	&1,467	&2,296	&3,592	&10,320	&6,159,463 \\
Num Search Results	&10	&20	&26	&31	&46	&2,095 \\
Search Results Size	&1,1691	&33,989	&49,222	&68,681	&135,533	&5,355,671 \\
\end{tabular}
\end{center}
\label{wikisum-stats}
\end{table}
  % Note: auto-generated by csv_to_table.py then manually added "," to numbers

Table \ref{wikisum-stats} describes overall properties of our WikiSum dataset.
Many articles have few citations, motivating our supplementation of the source
documents with web search results. On the other hand, citations when
available, tend to be of higher-quality.  When counting the total words
in the entire dataset, it is orders-of-magnitude larger than previous summarization
datasets.

To have consistent train/development/test data across corpus-comparison experiments,
 we restrict the articles to those with at least one crawlable citation.
We divide the articles roughly into 80/10/10 for train/development/test subsets,
resulting in 1865750, 233252, and 232998 examples respectively.
